\begin{textsample}{POS Dim 3 – rr – Score 11.00 – rr/rr\_inf\_4.txt}  \label{ex:f3_pos_157}
3.1.3 DIREITOS DE APRENDIZAGEM O currículo da Educação \textbf{Infantil} necessita articular a promoção de experiências e vivências em torno do conhecimento produzido pela \textbf{criança}, oportunizando o desenvolvimento de suas capacidades : como a memória, a percepção, a interação, a imaginação, o pensamento, a consciência, o autoconhecimento, a afetividade, o movimento, as diferentes linguagens.
Nesse sentido, o currículo Estadual de Roraima fundamenta -se com base no Artigo 9º das Diretrizes Curriculares Nacionais para Educação \textbf{Infantil} – DCNEIS ( 2009 ), revelando que o currículo da Educação \textbf{Infantil} deve promover experiências diversas que ampliem o contato entre a \textbf{criança} e as diferentes culturas, no qual é essencial ao seu desenvolvimento : - CONVIVER com outras \textbf{crianças} e \textbf{adultos}, em \textbf{pequenos} e grandes grupos, utilizando diferentes linguagens, ampliando o conhecimento de si e do outro, o respeito em relação à cultura e às diferenças entre as pessoas. - \textbf{BRINCAR} cotidianamente de diversas formas, em diferentes espaços e tempos, com diferentes parceiros ( \textbf{crianças} e \textbf{adultos} ), ampliando e diversificando seu acesso a produções culturais, seus conhecimentos, sua imaginação, sua criatividade, suas experiências emocionais, corporais, \textbf{sensoriais}, expressivas, cognitivas, sociais e relacionais. - PARTICIPAR ativamente, com \textbf{adultos} e outras \textbf{crianças}, tanto do planejamento da gestão da escola e das atividades propostas pelo educador quanto da realização das atividades da vida cotidiana, tais como a escolha das \textbf{brincadeiras}, dos materiais e dos ambientes, desenvolvendo diferentes linguagens e elaborando conhecimentos, decidindo e se posicionando. - EXPLORAR movimentos, \textbf{gestos}, \textbf{sons}, formas, \textbf{texturas}, cores, palavras, emoções, transformações, relacionamentos, histórias, objetos, elementos da natureza, na escola e fora dela, ampliando seus saberes sobre a cultura, em suas diversas modalidades : as artes, a escrita, a ciência e a tecnologia. - EXPRESSAR, como sujeito dialógico, criativo e sensível, suas necessidades, emoções, sentimentos, dúvidas, hipóteses, descobertas, opiniões, questionamentos, por meio de diferentes linguagens. - CONHECER -SE e construir sua identidade pessoal, social e cultural, constituindo uma imagem positiva de si e de seus grupos de pertencimento, nas diversas experiências de \textbf{cuidados}, interações, \textbf{brincadeiras} e linguagens vivenciadas na instituição escolar e em seu contexto familiar e comunitário.

% matched lemmas: adulto, brincadeira, brincar, criança, cuidado, gesto, infantil, pequeno, sensorial, som, textura
\end{textsample}
